\section{Discussion and Limitations}\label{sec:discussion}

The evaluation shows that \sysname{} improves over GraphCL on all three datasets, but the size of the improvement differs by dataset. The gain is small on StreamSpot and much larger on GraSec-IoT and Unicorn Wget. This suggests that topology-aware contrastive learning is most useful when malicious activity changes graph structure in ways that are captured by the topology-aware representation.

The dataset-level results also show that the value of supervision depends strongly on how well the labeled attacks represent the test behavior. The supervised GNN remains the strongest model on StreamSpot and GraSec-IoT. This is expected because it uses labeled attack examples during training, and these datasets provide attack patterns that are more directly learnable from the available labels. StreamSpot consists of controlled task-level flow graphs from repeated non-malicious browser activities and GraSec-IoT is designed as an IoT graph dataset for attack detection. In these settings, labeled attack examples can help a supervised GNN learn a useful decision boundary. In contrast, Unicorn is a system-level dataset based on a APT scenario. In this dataset non-malicious and malicious graphs can have similar activity characteristics. This makes the dataset harder for a label-dependent graph classifier. The classifier overfits to local patterns rather than learning a robust structural boundary. Therefore, although the supervised GNN achieves a non-trivial score on Unicorn Wget, it is less stable and weaker than TopoIDS, especially when recall and F1 are considered.


The supervised GNN remains the strongest model on StreamSpot and GraSec-IoT. This is expected because it uses labeled attack examples during training. However, the supervised GNN performs poorly on Unicorn Wget, especially in recall and F1. This result shows a practical risk of label-dependent intrusion detection: when labeled attack training data does not match the test behavior, a supervised model may miss many attacks.

The main value of \sysname{} is that it learns representations from non-malicious graphs and improves over a strong contrastive baseline. Its scoring design keeps representation learning and anomaly detection consistent by using the same contrastive embedding space for both training and scoring. This design does not require labeled attacks during representation learning, an extra classifier, or a larger model.

This work has some limitations. First, the gain on StreamSpot is small, which shows that topology-aware scoring does not always produce a large improvement. Second, \sysname{} performs graph-level detection. It predicts whether a whole graph is non-malicious or malicious, but it does not identify which nodes, edges, or subgraphs caused the alert. Future work should extend the method to support attack localization and alert explanation.